% TODO: Make a glossary
% All \newacronym and \newglossaryentry definitions live here


% Computer Resources =============================================
\newglossaryentry{daw}
{
    name={DAW},
    description={Digital Audio Workstation; software used for recording, editing, processing, and producing digital audio}
}

\newglossaryentry{io}
{
    name={I/O},
    plural={I/Os},
    description={Input/Output; the communication between a computer 
    system and external devices or software components}
}

\newglossaryentry{cpu}
{
    name={CPU},
    description={Central Processing Unit; the primary processor responsible for executing program instructions}
}

\newglossaryentry{ram}
{
    name={RAM},
    description={Random Access Memory; volatile memory used to store data and program state during execution}
}

\newglossaryentry{api}
{
    name={API},
    description={Application Programming Interface; a set of rules and functions that allows software components to communicate with one another}
}

\newglossaryentry{uml}
{
    name={UML},
    description={Unified Modeling Language; a standardized visual modeling language used to represent the 
    structure, behavior, and relationships of software systems through diagrams such as class diagrams 
    and sequence diagrams}
}

\newglossaryentry{dao}
{
    name={DAO},
    description={Data Access Object; a software design pattern that separates data storage and 
    retrieval operations from the application logic by providing an abstraction layer for 
    accessing data}
}

% Audio sampling and DSP ==============================================
\newglossaryentry{adc}
{
    name={ADC},
    description={Analog-to-Digital Converter; a device that converts a continuous analog
    signal into a discrete digital representation by sampling and
    quantizing its amplitude}
}

\newglossaryentry{dac}
{
    name={DAC},
    description={Digital-to-Analog Converter; a device that converts a discrete digital
    signal back into a continuous analog signal by reconstructing the
    waveform from quantized amplitude values}
}

\newglossaryentry{dsp}
{
    name={DSP},
    description={Digital Signal Processing; the use of digital processing to analyze, modify, and 
    synthesize signals, such as sound, by converting them into a digital format for 
    mathematical manipulation.}
}


\newglossaryentry{artifact}
{
    name={Artifact},
    description={Undesirable noises that occur during the recording, editing, or manipulation of audio}
}

\newglossaryentry{iir}
{
    name={IIR},
    description={Infinite Impulse Response; a digital filter where the impulse response does not 
    settle to zero but continues indefinitely}
}

\newglossaryentry{fir}
{
    name={FIR},
    description={Finite Impulse Response; a digital filter whose impulse response settles to zero after a certain amount of time}
}

\newglossaryentry{fft}
{
    name={FFT},
    description={Fast Fourier Transform; an efficient algorithm for computing the discrete Fourier transform}
}

\newglossaryentry{sf2}
{
    name={SF2},
    description={SoundFont 2; a file format for sample-based virtual instruments}
}

\newglossaryentry{wav}
{
    name={WAV},
    description={Waveform Audio File Format; a standard file format for storing uncompressed digital audio}
}

\newglossaryentry{riff}
{
    name={RIFF},
    description={Resource Interchange File Format; a container format used to store multimedia data, including WAV audio files in tagged chunks}
}

% Musical instruments ============================================
\newglossaryentry{sfz}{
    name={SFZ},
    description={an open, plain-text file format for defining
    sample-based instruments}
}

\newglossaryentry{fluidsynth}{
    name={FluidSynth},
    description={free open source software synthesizer which converts MIDI 
    note data into an audio signal using SoundFont technology without the need 
    for a SoundFont-compatible soundcard}
}

\newglossaryentry{synthesizer}
{
    name={Synthesizer},
    description={An electronic instrument or software system that generates audio signals using 
    oscillators/pre-recorded wavetables, filters, envelopes, and modulation sources to create sounds}
}

\newglossaryentry{oscillator}
{
    name={Oscillator},
    description={A signal generator that produces a periodic waveform, such as a sine, 
    square, triangle, or sawtooth wave, used as a fundamental sound source for synthesizers}
}

\newglossaryentry{unison}
{
    name={Unison},
    description={a technique where multiple voices or oscillators play the same note 
    simultaneously, often with slight variations in pitch, timing, or phase, known 
    for creating a thicker, wider, and more powerful sound}
}

\newglossaryentry{fourieranalysis}
{
    name={Fourier Analysis},
    description={A mathematical technique that decomposes a signal into a set of sinusoidal 
    components with different frequencies, amplitudes, and phases}
}

% =================================================================
\newglossaryentry{samplerate}
{
    name={Sample Rate},
    description={The number of samples taken from a continuous audio signal per second 
    during analog-to-digital conversion, measured in Hertz}
}

\newglossaryentry{midi}
{
    name={MIDI},
    description={Musical Instrument Digital Interface; a technical standard that allows electronic
    musical instruments, computers, and other devices to communicate with each other}
}


\newglossaryentry{midimessage}
{
    name={MIDI Message},
    description={A data packet defined by the MIDI protocol that communicates musical information, 
    such as note activation, note release, controller changes, and other performance parameters 
    between MIDI devices or software components}
}

\newglossaryentry{midichannel}
{
    name={MIDI Channel},
    description={A logical communication channel within the MIDI protocol that allows multiple 
    independent instruments or tracks to share a single MIDI connection, with up to 16 channels 
    available per MIDI port}
}

\newglossaryentry{midievent}
{
    name={MIDI Event},
    description={A specific action or message in the MIDI protocol which the MIDI device sends/receives
        in order to perform a task, depending on the details like the note number and velocity}
}

\newglossaryentry{bpm}
{
    name={BPM},
    description={Beats Per Minute; a unit used to specify the tempo of music}
}


\newglossaryentry{cent}
{
    name={Cent},
    description={A logarithmic unit of musical pitch interval where one semitone 
    is divided into 100 cents, and 1200 cents correspond to one octave}
}


\newglossaryentry{adsr}
{
    name={ADSR},
    description={Attack, Decay, Sustain, Release; the four stages of an envelope used to control the amplitude or other 
    parameters of a sound as long as the key in the MIDI controller is pressed}
}

\newglossaryentry{lfo}
{
    name={LFO},
    description={Low-Frequency Oscillator; an oscillator operating below the audible frequency range, 
    typically used to modulate synthesizers and effects, creating rhythmic pulses}
}

% =========== FILTER ==================
\newglossaryentry{filter}{
    name={Filter},
    description={frequency-dependent system that can amplify, pass, or attenuate 
    specific frequency ranges within the audio frequency range of 0 Hz to 20 kHz}
}

\newglossaryentry{biquad}{
    name={Biquad Filter},
    description={A recursive second-order digital filter used to implement various filter types
    efficiently}
}

\newglossaryentry{lpf}
{
    name={LPF},
    description={Low-Pass Filter; system that allows signals with a frequency lower 
    than a certain cutoff frequency to pass through while attenuating signals 
    with frequencies higher than that cutoff}
}


\newglossaryentry{hpf}
{
    name={HPF},
    description={High-Pass Filter; system that allows signals with a frequency higher 
than a certain cutoff frequency to pass through while attenuating signals 
with frequencies lower than that cutoff}
}


\newglossaryentry{bpf}
{
    name={BPF},
    description={Band-Pass Filter; system that passes frequencies within a certain 
range and attenuates frequencies outside that range}
}


\newglossaryentry{peaking}{
    name={Peaking Filter},
    description={A filter that boosts or cuts a band of frequencies
    centered on a cutoff frequency, with the effect tapering off
    on either side, controlled by a Q factor}
}

\newglossaryentry{shelf}{
    name={Shelf Filter},
    description={A filter that boosts or cuts all frequencies above
    (high-shelf) or below (low-shelf) a chosen cutoff frequency by
    a constant amount, leaving the rest of the spectrum unaffected}
}

\newglossaryentry{notch}{
    name={Notch Filter},
    description={A narrow-band filter that introduces a sharp attenuation within a
    specific frequency while leaving surrounding frequencies largely
    unaffected}
}

\newglossaryentry{svf}
{
    name={SVF},
    description={State-Variable Filter; a filter structure capable of simultaneously producing 
    low-pass, high-pass, band-pass, and other filter responses, ideal for real-time frequency sweeps}
}


\newglossaryentry{zdf}
{
    name={ZDF},
    description={Zero-Delay Feedback; a digital filter design technique that eliminates 
    the one-sample delay normally introduced in feedback loops, commonly using algebra}
}

\newglossaryentry{rms}
{
    name={RMS},
    description={Root Mean Square; a statistical measure commonly used to represent the 
    effective amplitude or power of an signal or wave}
}

\newglossaryentry{rt60}
{
    name={RT60},
    description={Reverberation Time; the duration it takes for sound to decay by 60 decibels 
    in a space after the sound source has stopped}
}

\newglossaryentry{gain}
{
    name={Gain},
    description={The amount of amplification or attenuation applied to an audio signal, typically expressed as a linear factor or in decibels}
}


\newglossaryentry{compression}
{
    name={Compression},
    description={An audio processing technique that reduces the dynamic range of a signal by attenuating portions of the signal that exceed a specified threshold}
}

% ========= Reverb ======================
\newglossaryentry{fdn}
{
    name={FDN},
    description={Feedback Delay Network; a reverberation structure consisting of multiple 
    interconnected delay lines with feedback}
}

\newglossaryentry{diffusion}
{
    name={Diffusion},
    description={the scattering of sound waves in many directions to 
    reduce echoes and make the sound feel more natural and spacious 
    in a room}
}

\newglossaryentry{delay}
{
    name={Delay},
    description={An audio effect or signal processing technique that introduces a time lag between the input signal and the output, commonly used to create echoes and spatial effects}
}

% =============== Threading ========================
\newglossaryentry{dag}
{
    name={DAG},
    description={Directed Acyclic Graph; a graph of nodes connected by directed
        edges containing no cycles, commonly used to represent a
        processing or signal-flow chain with a defined execution
        order}
}


\newglossaryentry{criticalpath}
{
    name={Critical Path},
    description={The longest path through the DAG, representing the sequence of tasks 
    that determines the maximum processing time}
}

\newglossaryentry{raii}
{
    name={RAII},
    description={A programming idiom in C++ where resources are acquired during 
    object creation and automatically released when the object goes out of scope}
}
